% !TEX program = xelatex
% Sources: R1ezzTa/MotionJPEG, remote commit 4c5a664, accessed 2026-10-24.
% Compile twice. Notes are hidden; uncomment show notes for speaking prompts.
\documentclass[aspectratio=169,11pt]{beamer}
\usetheme{Madrid}
\usepackage[UTF8,fontset=fandol]{ctex}
\usepackage{amsmath,booktabs,array,graphicx,tikz,listings,appendixnumberbeamer}
\usetikzlibrary{arrows.meta,positioning,calc,fit}
\definecolor{navy}{HTML}{262681}
\definecolor{teal}{HTML}{3333AC}
\definecolor{ink}{HTML}{111111}
\definecolor{muted}{HTML}{666666}
\definecolor{pale}{HTML}{E8E9F1}
\setbeamercolor{normal text}{fg=ink,bg=white}
\setbeamercolor{structure}{fg=teal}
\setbeamercolor{frametitle}{fg=white,bg=teal}
\setbeamercolor{title}{fg=white,bg=teal}
\setbeamercolor{block title}{fg=white,bg=teal}
\setbeamercolor{block body}{fg=ink,bg=pale}
\setbeamerfont{frametitle}{size=\large,series=\mdseries}
\setbeamerfont{title}{size=\Large,series=\mdseries}
\setbeamerfont{subtitle}{size=\normalsize}
\setbeamersize{text margin left=9mm,text margin right=9mm}
\setbeamertemplate{navigation symbols}{}
\setbeamertemplate{items}[ball]

\setbeamercolor{author in head/foot}{fg=white,bg=navy}
\setbeamercolor{title in head/foot}{fg=white,bg=teal}
\setbeamercolor{date in head/foot}{fg=white,bg=teal!80!black}
\setbeamerfont{author in head/foot}{size=\fontsize{5.5}{6.5}\selectfont}
\setbeamerfont{title in head/foot}{size=\fontsize{5.5}{6.5}\selectfont}
\setbeamerfont{date in head/foot}{size=\fontsize{5.5}{6.5}\selectfont}
\setbeamertemplate{footline}{%
  \leavevmode\hbox{%
  \begin{beamercolorbox}[wd=.20\paperwidth,ht=2.5ex,dp=1.1ex,center]{author in head/foot}%
    \usebeamerfont{author in head/foot}\insertshortauthor
  \end{beamercolorbox}%
  \begin{beamercolorbox}[wd=.53\paperwidth,ht=2.5ex,dp=1.1ex,center]{title in head/foot}%
    \usebeamerfont{title in head/foot}\hypersetup{linkcolor=white}\insertshorttitle
  \end{beamercolorbox}%
  \begin{beamercolorbox}[wd=.27\paperwidth,ht=2.5ex,dp=1.1ex,right]{date in head/foot}%
    \usebeamerfont{date in head/foot}\insertshortinstitute\hspace{2mm}%
    \ifnum\value{framenumber}>\inserttotalframenumber\relax
      附录~\number\numexpr\value{framenumber}-\inserttotalframenumber\relax
    \else\insertframenumber\,/\,\inserttotalframenumber\fi\hspace{2mm}
  \end{beamercolorbox}}%
}

\setlength{\parskip}{4pt}
\setlength{\tabcolsep}{5pt}
\renewcommand{\arraystretch}{1.23}
\newcommand{\key}[1]{{\color{teal}\bfseries #1}}
\newcommand{\source}[1]{\par\vfill{\fontsize{6.8}{8}\selectfont\color{muted}依据：#1\par}}
\newcommand{\smallnote}[1]{\par{\footnotesize\color{muted}#1\par}}
\newcommand{\repo}[2]{\href{https://github.com/R1ezzTa/MotionJPEG/blob/4c5a664288e3a78fe1c78c0bbdc8db61b8cb6bdf/#1}{#2}}
\tikzset{box/.style={draw=teal!65,fill=pale,minimum height=10mm,text width=25mm,align=center,font=\small,inner sep=3pt}, arr/.style={-{Stealth[length=2mm]},thick,draw=navy}, lab/.style={font=\scriptsize,text=muted,align=center}}
\lstset{language=Verilog,basicstyle=\ttfamily\footnotesize,keywordstyle=\color{teal}\bfseries,commentstyle=\color{muted},backgroundcolor=\color{pale!65},frame=none,columns=fullflexible,keepspaces=true,showstringspaces=false,breaklines=true,aboveskip=6pt,belowskip=6pt}
\hypersetup{colorlinks=true,urlcolor=teal,linkcolor=navy,pdftitle={MotionJPEG：FPGA实时图像采集与JPEG编码},pdfauthor={任子桐}}
\title[MotionJPEG：FPGA 实时图像采集与 JPEG 编码]{MotionJPEG}
\subtitle{无 DDR3 的 FPGA 实时图像采集与 JPEG 编码}
\author{任子桐}
\institute[浙江大学]{浙江大学\\硬件描述语言课程项目}
\date{2026 年 10 月}
\begin{document}
\begin{frame}{MotionJPEG：FPGA 实时图像采集与编码}
\begin{beamercolorbox}[wd=\textwidth,sep=4mm,rounded=true,shadow=true]{title}
\centering\Large 在片内完成图像处理与 JPEG 编码\par
\vspace{2mm}\normalsize OV5640 摄像头输入，USB 输出至 PC
\end{beamercolorbox}
\vspace{5mm}
\begin{columns}[T,onlytextwidth]
\column{.53\textwidth}
\key{项目功能}\par
\vspace{1mm}
彩色、灰度、二值化、Sobel 四种模式\par
JPEG 质量档位 Q60 / Q75 / Q85\par
每帧独立编码，PC 接收、解码与显示
\column{.41\textwidth}
\key{设计约束}\par
\vspace{1mm}
Artix-7 XC7A35T\par
不使用 DDR3 整帧缓存\par
采用行缓存、块缓存与流式处理
\end{columns}
\vfill
{\small 任子桐\hfill 硬件描述语言课程项目 · 浙江大学}
\end{frame}

\begin{frame}{系统架构：数据跨时钟域连续流动}
\centering
\begin{tikzpicture}[node distance=9mm and 19mm]
\node[box,text width=30mm] (a) {摄像头采集\\OV5640 / DVP};
\node[box,text width=30mm,right=of a] (b) {图像处理\\颜色重建与降噪};
\node[box,text width=30mm,right=of b] (c) {JPEG 编码\\分块、变换与压缩};
\node[box,text width=30mm,below=12mm of c] (d) {帧封装\\MBLK 数据包};
\node[box,text width=30mm,left=of d] (e) {USB 发送\\FT245 接口};
\node[box,text width=30mm,left=of e] (f) {PC 接收\\解码与显示};
\draw[arr] (a)--node[above,lab]{异步 FIFO}(b);
\draw[arr] (b)--(c);\draw[arr] (c)--(d);
\draw[arr] (d)--node[above,lab]{异步 FIFO}(e);\draw[arr] (e)--(f);
\node[lab,above=4mm of a]{摄像头时钟域};
\node[lab,above=4mm of c]{100 MHz 核心域};
\node[lab,below=3mm of e]{60 MHz USB 域};
\end{tikzpicture}
\vspace{4mm}
\begin{block}{两个接口问题}
\small
\key{跨域：}FIFO 的 Gray 指针经过两级同步，数据通过双时钟存储体传递。\par
\key{停顿：}反压时同步保持数据、有效位和帧标记，避免像素与帧边界错位。
\end{block}
\end{frame}

\begin{frame}{关键实现：存储并行化，运算分级寄存}
\begin{columns}[T,onlytextwidth]
\column{.47\textwidth}
\key{1. 像素与系数怎样存}\par
\small
光栅输入按存储体分开，支持双样本访问。\par
Zigzag 通过分布式 RAM 重排系数。
\vspace{3mm}
\begin{tikzpicture}[node distance=7mm]
\node[box,text width=18mm,font=\footnotesize] (a){光栅像素};
\node[box,text width=18mm,font=\footnotesize,right=of a] (b){分块 RAM};
\node[box,text width=18mm,font=\footnotesize,below=6mm of b] (c){8$\times$8 块};
\draw[arr](a)--(b);\draw[arr](b)--(c);
\end{tikzpicture}
\column{.49\textwidth}
\key{2. 组合路径怎样拆}\par
\small
\texttt{formatter\_start} 延迟一拍。\par
DCT 用局部行驱动寄存器分散扇出。\par
量化用倒数乘法加修正替代直接除法。
\vspace{3mm}
\begin{tikzpicture}[node distance=4mm]
\node[box,text width=16mm,font=\footnotesize] (a){倒数查表};
\node[box,text width=16mm,font=\footnotesize,right=of a] (b){乘法};
\node[box,text width=16mm,font=\footnotesize,right=of b] (c){比较修正};
\draw[arr](a)--(b);\draw[arr](b)--(c);
\node[lab,below=3mm of b]{各级寄存，控制信号随数据对齐};
\end{tikzpicture}
\end{columns}
\vspace{3mm}
\begin{block}{3. 系数怎样压缩}
\small DC 差分表示相邻块变化；AC 零游程表示连续零；EOB 结束块内剩余零。
\end{block}
\end{frame}

\begin{frame}{验证结果：功能通过，资源满足器件容量}
\begin{columns}[T,onlytextwidth]
\column{.47\textwidth}
\key{功能验收基线}\par
{\fontsize{30}{34}\selectfont\color{teal}1080p30}\par
\small 多项测试累计 \key{3480 帧}独立解码通过\par
摄像头 FIFO 最大溢出计数：\key{0}\par
\vspace{2mm}
Sobel 连续运行 60 s\par
1800 帧，平均 30.0019 fps
\column{.46\textwidth}
\key{实现资源}\par
\small
\begin{tabular}{@{}lrr@{}}
\toprule
资源 & 使用量 & 占比\\\midrule
LUT & 15,584 & 74.92\%\\
FF & 19,533 & 46.95\%\\
BRAM36 & 36 & 72.00\%\\
DSP & 46 & 51.11\%\\\bottomrule
\end{tabular}\par
\vspace{2mm}
WNS：$+0.057$ ns\par
WHS：$+0.023$ ns
\end{columns}
\vspace{3mm}
\begin{block}{验证结论}
\small 已完成无 DDR3 的实时编码链路；后续精简版本通过回归与综合，仍需重新验收。
\end{block}
\source{功能、资源与时序对应 full\_20260929\_170515。3480 帧为多项测试累计，BRAM 按 36 Kb 等效计。}
\end{frame}

\begin{frame}{项目分工}
\small
\renewcommand{\arraystretch}{1.65}
\begin{tabular}{@{}>{\raggedright\arraybackslash}p{2.25cm}>{\raggedright\arraybackslash}p{6.8cm}p{2.5cm}@{}}
\toprule
\key{工作方向} & \key{负责内容} & \key{成员}\\\midrule
采集与处理 & 摄像头初始化、采集、颜色重建与预处理 & 待填写\\
JPEG 编码 & 分块缓存、DCT、量化与符号编码 & 待填写\\
传输与上位机 & USB 封包、PC 接收与解码显示 & 待填写\\
集成与验证 & 顶层联调、仿真回归及验收整理 & 待填写\\\bottomrule
\end{tabular}
\vspace{5mm}
\smallnote{目前按工作方向列出，成员及实际职责待确认后填写。}
\vfill
\centering\large\color{teal}谢谢各位老师和同学
\end{frame}
\end{document}
